\section{Empirical Strategy}
\label{sec:empirical}

The analysis has two stages: we first estimate the overall cost of urgency by comparing ordinary and urgent purchases, then isolate the sanction channel by comparing litigated and administrative urgent purchases.

\subsection{Identification}

Identification requires that, conditional on fixed effects, urgency is uncorrelated with unobserved determinants of procurement outcomes. Three institutional features support this. First, court orders originate from patients' health needs and legal claims---not from the procurement process. Brazilian public administration operates under a principle of impersonality: tender outcomes depend on item specifications and market conditions, not on who requested the purchase. Second, court orders are unpredictable from the buyer's perspective. Only \BPstockShortageShare{} of claims arise from stock shortages or service failures, so litigation is poorly correlated with PBU-level mismanagement. Third, the urgency regime is determined externally---by judges or by the SES/SP's scientific committee---not by the purchasing unit. While PBUs could in principle misclassify purchases, the legal requirement to reference court orders in tender notices, combined with our REGEX-based classification of official texts, limits this concern. An honest-DiD event study around each item's first court order, computed via the imputation estimator of \citet{borusyak2024revisiting} (Online Appendix, Figure~\ref{fig:event_study}), exhibits pre-period coefficients statistically and economically indistinguishable from zero---supporting parallel trends---and post-period coefficients that rise monotonically from $+\BPbjsETzeroPct$ at $t=0$ to $+\BPbjsETfivePct$ at $t=+5$. This pattern rules out the reverse-causality story in which high prices attract lawsuits.

\paragraph{Selection into the administrative channel.} The under-the-gun comparison rests on the institutional fact that administrative and litigated urgent purchases share planning constraints and differ only in penalty exposure. Cases enter the administrative channel rather than the courts when the patient (or the patient's caregiver, often via SUS social workers) requests the medication directly through an SES/SP intake, and the request is approved by a scientific committee using cost-effectiveness criteria. Cases that bypass the administrative channel either (a) were not submitted to it because the patient went straight to litigation, or (b) were rejected by the committee. The committee thus selects on observed and unobserved item characteristics (essentiality, drug class, expected cost, evidence of clinical benefit). The direction of the resulting bias on our UTG coefficient is theoretically ambiguous: if the committee admits cost-effective items and rejects expensive outliers, our administrative subsample under-states the average cost of urgent procurement and our UTG gap is biased upward; if the committee admits items that are simpler to source and rejects clinically complex ones, our administrative subsample under-states the difficulty of the purchase and our UTG gap is biased downward. We do not attempt to sign this bias and instead frame the UTG estimate as a within-regime decomposition rather than a structural counterfactual against a hypothetical sanction-free regime; the latter object is unidentified in our setting.

\paragraph{Within-item balance.} Item fixed effects absorb time-invariant item characteristics but not within-item composition shifts: if severity, dose, or packaging systematically differ across purchase types, our estimates conflate procurement-channel effects with case mix. Within-item balance on pre-determined covariates (Online Appendix, Table~\ref{tab:balance_within}) supports the identifying assumption on observables: SUS-formulary status, time period, and buyer size are jointly balanced ($p > \BPbalPMin$ for each). The two dimensions that differ within item---log reference price and log quantity---are precisely the outcomes the paper attributes to the sanction mechanism (officials set looser reference prices and place smaller orders under judicial pressure); their imbalance is a \emph{feature} of the treatment, not a confounder. Procurement modality (sealed-bid vs.\ reverse auction) differs by \BPbalModalityPP, a gap small relative to the \BPutgPctRange{} UTG premium and absorbed by year-month fixed effects.

\paragraph{Other threats.} Three further threats warrant discussion. \emph{First}, court orders originate from patients' health needs and legal claims, not from the procurement process. Brazilian public administration operates under a principle of impersonality: tender outcomes depend on item specifications and market conditions, not on who requested the purchase. Only \BPstockShortageShare{} of claims arise from stock shortages or service failures, so litigation is poorly correlated with PBU-level mismanagement. \emph{Second}, an honest-DiD event study around each item's first court order, computed via the imputation estimator of \citet{borusyak2024revisiting} (Online Appendix, Figure~\ref{fig:event_study}), exhibits pre-period coefficients economically and statistically indistinguishable from zero---supporting parallel trends---and post-period coefficients that rise from $+\BPbjsETzeroPct$ at $t=0$ to $+\BPbjsETfivePct$ at $t=+5$. The dynamic event-study evidence rules out the reverse-causality story in which high prices attract lawsuits and complements the cross-sectional fixed-effects design. \emph{Third}, regex-based classification of tender notices may contain asymmetric misclassification between administrative and litigated requests; the classifier's design and hand-labeling protocol are documented in Online Appendix \S A.7. Together with the never-litigated placebo (Section~\ref{sec:placebo}), the design narrows the residual identification space without eliminating it. Throughout the paper we frame the UTG estimate as a descriptive decomposition rather than a structural parameter; this is the central interpretive claim, and we honor it by reporting an explicit channel decomposition rather than a single counterfactual price.

\subsection{Enforcement Costs: Ordinary vs.\ Urgent}

We estimate the following baseline specification:
\begin{equation}
\label{eq:baseline}
y_{i,g,t} = \beta \cdot \text{Urgent}_{i,g,t} + \gamma_g + \lambda_t + \delta_b + \varepsilon_{i,g,t},
\end{equation}
where $y_{i,g,t}$ is the outcome for purchase order $i$ of item $g$ at time $t$; $\gamma_g$ are item fixed effects; $\lambda_t$ are time fixed effects; and $\delta_b$ are PBU fixed effects. Standard errors are clustered at the PBU level. We report four specifications with progressively richer fixed effects, from item FE only through item + year-month + PBU FE. Our preferred specification (item + year + PBU FE) absorbs time-invariant item characteristics, common annual trends, and buyer-specific heterogeneity, while retaining sufficient within-cell variation for precise estimation.

For negotiated prices and number of firms, we additionally estimate a specification that partials out log quantity:
\begin{equation}
\label{eq:direct}
y_{i,g,t} = \beta \cdot \text{Urgent}_{i,g,t} + \alpha \cdot \ln Q_{i,g,t} + \gamma_g + \lambda_t + \delta_b + \varepsilon_{i,g,t},
\end{equation}
Because log quantity is itself an outcome of urgency, \eqref{eq:direct} does not identify a structural direct effect \citep{acharya2016explaining}; we therefore read the contrast between \eqref{eq:baseline} and \eqref{eq:direct} as a descriptive mediation decomposition, quantifying how much of the total premium co-moves with the quantity margin that urgency itself induces. Section~\ref{sec:main_results} develops this interpretation.

\subsection{The ``Under the Gun'' Effect: Litigated vs.\ Administrative}

To isolate the effect of judicial sanctions, we restrict the sample to urgent purchases only and estimate:
\begin{equation}
\label{eq:utg}
\ln(\text{Price}_{i,g,t}) = \beta \cdot \text{Admin}_{i,g,t} + \gamma_g + \lambda_t + \delta_b + \varepsilon_{i,g,t},
\end{equation}
where $\text{Admin}$ equals 1 for administrative purchases and 0 for litigated. Since both types face identical planning constraints, $\beta$ isolates the cost attributable to the sanction regime.
